\par\textbf{Procedimiento completo para comparar tu solución.} El giro tiene aceleración aunque la rapidez sea constante. Equilibrio vertical fija $n=1/\cos\phi$; la componente horizontal da el radio. Se comprueban por separado sustentación y resistencia al aumentar la carga.

\begin{enumerate}

\item Factor de carga de giro.

\[\fitmath{n=\frac{1}{\cos\left(\frac{30\cdot \pi}{180}\right)}\approx 1.154701\,}\]

\item Radio por aceleración centrípeta.

\[\fitmath{r=\frac{{55}^{2}}{9.80665\cdot \tan\left(\frac{30\cdot \pi}{180}\right)}\approx 534.2756\,\mathrm{m}}\]

\item Fuerza de referencia a 55 m/s.

\[\fitmath{\mathrm{qS}=\frac{0.95\cdot {55}^{2}\cdot 24.36}{2}\approx 35002.27\,\mathrm{N}}\]

\item Sustentación requerida en giro.

\[\fitmath{C_L=\frac{n\cdot 40000}{\mathrm{qS}}}\]

\[\fitmath{C_L=\frac{(1.154700538)\cdot 40000}{(35002.275)}\approx 1.319572\,}\]

\item Resistencia del mismo estado de giro.

\[\fitmath{D=\mathrm{qS}\cdot \left(0.028+0.0519571\cdot {C_L}^{2}\right)}\]

\[\fitmath{D=(35002.275)\cdot \left(0.028+0.0519571\cdot {(1.319571986)}^{2}\right)\approx 4146.767\,\mathrm{N}}\]

\item Margen de empuje; comparar además CL con 1.4.

\[\fitmath{\mathrm{margin}=4000-D}\]

\[\fitmath{\mathrm{margin}=4000-(4146.766817)\approx -146.7668\,\mathrm{N}}\]

\end{enumerate}
